% Generated draft: requires mathematical and source review.
\documentclass[11pt]{article}
\usepackage[margin=25mm]{geometry}
\usepackage{fontspec}
\setmainfont{Latin Modern Roman}
\usepackage{amsmath,amssymb,mathtools,unicode-math}
\setmathfont{Latin Modern Math}
\usepackage{enumitem,needspace,xurl,hyperref}
\setlength{\parindent}{0pt}
\setlength{\parskip}{5pt}
\setlength{\emergencystretch}{3em}
\newcommand{\N}{\mathbb{N}}
\newcommand{\Z}{\mathbb{Z}}
\newcommand{\Q}{\mathbb{Q}}
\newcommand{\R}{\mathbb{R}}
\newcommand{\C}{\mathbb{C}}
\newcommand{\F}{\mathbb{F}}
\newcommand{\A}{\mathbb{A}}
\newcommand{\PP}{\mathbb{P}}
\newcommand{\E}{\mathbb{E}}
\newcommand{\eps}{\varepsilon}
\newcommand{\dd}{\,\mathrm{d}}
\begin{document}
\section*{UnsolvedMath ID: 30000302}
\textbf{Source problem code:} catalogue entry
\subsection{Definitions and mathematical statement}
For observed DNA or protein sequences $x=(x_1,\ldots,x_m)$ and $y=(y_1,\ldots,y_n)$ over their respective finite alphabet, the source asks whether their similarity is attributable to chance or instead indicates common ancestry or function. The supplied assessment identifies this as an introductory local-alignment significance question, not a fully specified statistical theorem.

A precise formulation requires a null law $P_0$ for unrelated random sequences $(X,Y)$ of the observed lengths and a similarity statistic. For a local-alignment score one may write
\[
 S(X,Y)=\max_{A\text{ admissible local alignment}}\operatorname{score}(A;X,Y),
\]
where the substitution scores and any gap penalties must be fixed. Given $s_{\mathrm{obs}}=S(x,y)$, a null tail probability would be
\[
 p_0=P_0\{S(X,Y)\geq s_{\mathrm{obs}}\}.
\]
The maximum must be included in the null calibration if it was used to select the observed alignment. A level-$\alpha$ rejection rule needs a threshold with $P_0\{S\geq c_\alpha\}\leq\alpha$, or a corresponding calibrated rule for a specified composite null.

These formulas expose the statistical task without selecting an independent-letter, Markov, composition-preserving or evolutionary model. A significant score only supplies evidence relative to that null; it is not a posterior probability of ancestry or a proof of shared function. Assessing common ancestry and common function requires suitable, potentially different alternatives and power/error criteria. Neither the null, alternatives, scoring system, lengths/asymptotic regime nor significance threshold is supplied, so no universally valid biological decision or particular extreme-value theorem is asserted.

The literature triage checked 21 August 2026 says the general model is too underspecified to identify a solved theorem. The registry nevertheless labels the item open. These are preserved as contextual metadata, not evidence that every particular sequence-alignment significance problem is unsolved.
\subsection{Short English statement}
Assess whether an observed local sequence similarity is unusually large under a specified chance model. The score, null and alternative models and calibration are missing. Statistical significance alone does not establish ancestry or shared function, and the supplied assessment treats this as an introductory research question.
\subsection{Sources}
\begin{itemize}
\item Local UnsolvedMath registry, permanent ID 30000302, exact source code catalogue entry; \emph{Statistical Significance of Biological Sequence Similarity}. Proposed year 2005, status open, exact supplied sentence.
\item Supplied report citation: \emph{Mathematical Population Genetics (2006)}, Oberwolfach Reports, DOI 10.4171/owr/2005/40, pp. 2241--2304, \url{https://doi.org/10.4171/owr/2005/40}. Title year differs from DOI/proposed year.
\item Supplied literature triage checked 2026-08-21, citing only ``original Oberwolfach report'', \url{https://doi.org/10.4171/owr/2005/40}. It identifies an introductory local-alignment significance question with no sufficiently specified general statistical model. The dated OWR-derived curation requests independent verification of time-sensitive status and cited claims; no separate precise theorem reference is supplied.
\end{itemize}
\end{document}
